%%
\documentclass[article,screen,12pt, nonacm]{acmart}
\let\Bbbk\relax % Fix for amssymb clash 
\usepackage{standalone}
\usepackage{vmlmacros}
\AtBeginDocument{%
  \providecommand\BibTeX{{%
    \normalfont B\kern-0.5em{\scshape i\kern-0.25em b}\kern-0.8em\TeX}}}
\usepackage{outlines}
%~\usepackage{natbib}
\setlength{\headheight}{14.0pt}
\setlength{\footskip}{13.3pt}
\citestyle{acmauthoryear}
\begin{document}
\title{An Alternative to Pattern Matching, Inspired by Verse}
\author{Roger Burtonpatel}
\email{roger.burtonpatel@tufts.edu}
\affiliation{%
  \institution{Tufts University}
  \streetaddress{419 Boston Ave}
  \city{Medford}
  \state{Massachusetts}
  \country{USA}
  \postcode{02155}
}
\bibliographystyle{plainnat}
% \renewcommand\refname{}
% ~\setcitestyle{sort&compress}


%~\author{Norman Ramsey}
%~\email{nr@cs.tufts.edu}
%~\affiliation{%
%~\institution{Tufts University}
%~\streetaddress{177 College Ave}
%~\city{Medford}
%~\state{Massachusetts}
%~\country{USA}
%~\postcode{02155}
% }

%~\author{Milod Kazerounian}
%~\email{milod.mazerouniantufts.edu}
%~\affiliation{%
%~\institution{Tufts University}
%~\streetaddress{177 College Ave}
%~\city{Medford}
%~\state{Massachusetts}
%~\country{USA}
%~\postcode{02155}
% }

\renewcommand{\shortauthors}{Burtonpatel}

\rab{fix odd page indentation}
% Pretty dece: have Steve read. 
\begin{abstract}
    Pattern matching % S
    appeals % V
    to functional programmers for its expressiveness and good cost model, 
    but 
    it % S
    fails to express % V
    certain computations without verbosity. Verbosity % S 
    can be reduced % V
    by using extended forms of pattern matching,
    but extensions % (S = the problem)
    are % V
    not standardized, and making certain extensions like or-patterns and pattern
    guards cooperate in a large-scale programming language is a significant
    implementation challenge that has not seen success until recent years. 
    
    % the most desirable combination of extensions could not
    % until recently found in any one popular programming language.
    
    Alternatively to pattern matching, computations can be expressed succinctly
    by using equations, as has been demonstrated by the Verse programming
    language. But such computations may have time and space costs that can be
    hard to predict. 
    % But Verse's cost model % S-sub
    % is % V-sub
    % a challenge because of backtracking and multiple
    % results. 
    As a compromise, 
    we % S
    ~propose % V
    a new language,~\VMinus, % object 
    which~uses equations in a limited way that makes time and space costs easy
    to predict. 
    % to succinctly express computation while
    % enjoying a friendlier cost model by prohibiting backtracking and multiple
    % results. 
    In
    comparative examples,
    \VMinus~ % S 
    expresses % V
    computations as succinctly as pattern matching with popular extensions, and just like pattern matching, it % S
    can be compiled % V 
    to a decision tree. 
    
    % Evaluation
    To evaluate % v
    the succinctness of \VMinus,
    we compare % V
    it with a novel pattern-matching language, \PPlus. \PPlus
    demonstrates \rab{want better verb} % V
    three popular pattern matching extensions 
    in unison: side conditions, or-patterns, and pattern guards. 
    We evaluate % V 
    the expressiveness of both languages using four smaller contrived examples 
    and a larger example from \_. 

    An implementation of~\VMinus, \PPlus, and the match compiler can be found
    at~\url{https://github.com/rogerburtonpatel/vml}.
    % Question: implied that we've implemented D here, or not? 
    %%%%%%% Old, with "We show":
    % We % S
    % ~show % V
    %  that~\VMinus's equations capture the desirable properties of
    % pattern matching with popular extensions, 
    % and we % S-sub
    % ~also
    %  show % V 
    % that~\VMinus~can be compiled to a decision tree and therefore has
    % the same desirable cost model as pattern matching. Finally, 
    % we % S 
    % ~have % V 
    % implemented~\VMinus in Standard ML, alongside an algorithm for compiling it
    % to a decision tree. 
\end{abstract}

\maketitle

\section{Introduction}
% Subjects: Pattern matching and Equations 

Perhaps the most beloved tool among functional programmers for examining and
deconstructing data is pattern matching. 
% It does so implicitly by matching constructed data
% directly against a number of possible forms.
% Maybe combine? 
Pattern matching is also an established and well-researched topic
\citep{wadler1987views, macqueen1985tree, burton1993pattern, palao1996new,
maranget, bpc}. It is appreciated by programmers and researchers alike for two
main reasons: It enables~\it{implicit} data deconstruction, and it has a
desirable cost model. Specifically (regarding the latter), pattern matching can
be compiled to a~\it{decision tree}, a data structure that enforces linear
runtime performance by guaranteeing no part of the data will be examined more
than once.~\citep{maranget}

However, pattern matching cannot express certain common computations succinctly,
forcing programmers who wish to express these computations to duplicate code,
nest~\it{case} expressions, and create multiple points of truth. To mitigate
this, designers of popular programming languages have introduced~\it{extensions}
to pattern matching (Section~\ref{extensions}). 

Extensions strengthen pattern matching, but they are not standardized, so each
popular programming language with pattern matching features its own unique suite
of extensions. Extensions are subject to the discretion of the individual
language designer, not ubiquitous. Rather than continuing to extend pattern
matching~\it{ad hoc}, a worthwhile goal could be to find an alternative that
doesn't need extensions. A tempting possibility was introduced last year by the
programming language Verse~\citep{verse}. In Verse, a programmer can deconstruct
data using a different tool the language offers:~\it{equations}. Equations are
expressive and flexible, and it appears that they can express everything that
pattern matching can, including with popular extensions. 

But a full implementation of Verse is complicated, cost-wise. Verse is a
functional logic programming language, and expressions can backtrack at runtime
and return multiple results, both of which are hard to predict in their costs. 

% A worthwhile goal could be to harmonize the expressive quality of Verse's
% equations with the decision tree property of patterns.

In this paper, we~show that the expressive quality of Verse's equations and
the decision-tree property of patterns can be combined in a single language.
Since the language is a streamlined adaptation of Verse with a reduced feature
set, we~call it~\VMinus (“V minus”). For comparison, we present~\PPlus, a novel
pattern-matching language with three popular extensions. We compare~\VMinus
with~\PPlus in several small examples and a large example which is \padlexample. 


\rab{todo flow}

Our contributions are: 
We~have formalized~\PPlus and \VMinus with a big-step operational semantics
(Section~\ref{vminus}), we~have formalized decision trees into a core
language~\D (“D”) with a big-step operational semantics (Section~\ref{d}), 
% we~have formalized pattern matching in a core language~\PPlus (“P plus”) with a
% big-step operational semantics (Section~\ref{pplus}),
we~have formalized a translation from~\VMinus to~\D (Sections~\ref{vminustod}),
and we~have implemented both languages in Standard ML. 



\section{Pattern Matching and its Extensions}
\label{pmandextensions}
    \textbf{\rab{GOAL: Introduce pattern matching in a way that's neither
    dismissive of the interesting problem nor insulting to the audience's
    intelligence and knowledge of functional programming.}}


    Pattern matching is a battle-tested and beloved features in declarative and
    functional programming, but it often requires extensions to simplify cases
    that are otherwise troublesome. 
    
    
    Specifically, in sophisticated programs that involve complex
    decision-making, pattern matching without extensions tends to duplicate
    code. Consider this example of "simple" pattern matching (without
    extensions): 

    \rab{Example; will discuss which to use. Adaptation of 'cellar' example
    seems good}.

    \rab{Example says along the lines of: "Without extensions to pattern
    matching, expressing the algorithm looks like this: \dots}. 

    The code gets the job done, but it is cumbersome to write and read because
    of the large amount of duplication, which arises from the limitations of
    what can be expressed in simple pattern matching. 
    
    Furthermore, if the spec changes, will the programmer be shrewd enough to
    find and modify all the duplicated parts of the program to match the change?
    Such an example demands an improvement. Indeed, we can do better, by 
    augmenting pattern matching with \textit{extensions}. 


    \subsection{Popular extensions to pattern matching}
    \label{extensions}

        % Extensions to pattern matching simplify cases that are otherwise
    % troublesome. In this section, we~illustrate several such cases. The three
    % extensions we~describe are those commonly found in the literature and
    % implemented in compilers: side conditions, pattern guards, and or-patterns. 
    
    % To denote pattern matching~\it{without} extensions, we~coin the
    % term~\it{simple pattern matching}. In bare pattern matching, a pattern has
    % one of two syntactic forms: a name or an application (of a value constructor
    % to zero or more patterns). 


    \rab{TODO: Format this paragraph so it presents information as clearly as possible.}
    Duplicated code, especially in a complex case expression, can make code
    difficult to maintain and more bug-prone. To address this, many large-scale
    functional programming languages have added extensions to pattern matching,
    most commonly \textit{side conditions}(cite), which extend a pattern match with a
    boolean guard, \textit{or-patterns}(cite), which allow multiple branches in a case
    expression to share a right-hand side, and \textit{pattern guards}(cite), which
    allow multiple sequential pattern-matches and computations to guard a
    right-hand side. 

    \rab{Side condition example}
    Side conditions do \dots which eliminate the need to \dots, 
    illustrated in this small example of measuring shapes: 

        \rab{The examples are funny, but they should be more serious for this 
        submission.}

        \rab{TODO: merge figures, make them appear in the right place}
        \begin{figure}[ht]
        \begin{verbatim}
            type standard_shape = Square of float 
                    | Triangle of float * float
                    | Trapezoid of float * float * float
                    

            let exclaimTall sh =
            match sh with 
              | Square s -> if s > 100.0 
                            then "Wow! That's a sizeable square!"
                            else "Your shape is small." 
              | Triangle (_, h) -> 
                            if h > 100.0 
                            then "Goodness! Towering triangle!"
                            else "Your shape is small." 
              | Trapezoid (_, _, h) -> 
                            if h > 100.0
                            then "Zounds! Tremendous trapezoid!"
                            else "Your shape is small." 
            \end{verbatim}    
        \caption{An invented function~\tt{exclaimTall} in OCaml combines pattern
        matching with an~\tt{if} expression, and is not very pretty.}   
        \Description{An implementation of a function exclaimTall in OCaml, which
        uses pattern matching and an if statement.}
        \label{fig:ifexclaimtall}
    \end{figure}


       \begin{figure}[]
            \begin{verbatim}
                let exclaimTall sh =
                  match sh with 
                    | Square s when s > 100.0 ->
                            "Wow! That's a sizeable square!"
                    | Triangle (_, h) when h > 100.0 ->
                            "Goodness! Towering triangle!"
                    | Trapezoid (_, _, h) when h > 100.0 -> 
                            "Zounds! Tremendous trapezoid!"
                    | _ ->  "Your shape is small." 
                \end{verbatim}
            \caption{With a side condition,~\tt{exclaimTall} in OCaml becomes
            simpler and more adherent to the Nice~Properties.} 
            \Description{An implementation of a function exclaimTall in OCaml,
            which uses pattern matching and a side condition.}
            \label{fig:whenexclaimtall}
        \end{figure}


        \begin{figure}[ht]
        \begin{verbatim}

          let tripleLookup (rho : finitemap) (x : int) =
              match lookup rho x with Some w -> 
                (match lookup rho w with Some y -> 
                  (match lookup rho y with Some z -> z
                  | _ -> handleFailure x)
                | _ -> handleFailure x)
              | _ -> handleFailure x
            \end{verbatim}
        \caption{\tt{tripleLookup} in OCaml with bare pattern matching breaks
         the Nice~Property of~\hyperref[p2]{Single copy}: avoiding duplicated
         code.} 
                    
        \Description{An implementation of the tripleLookup with three levels of
        nested pattern matching.}
        \label{fig:pmtriplelookup}
    \end{figure}

     \begin{figure}
        \begin{center}
        \begin{verbatim}
                        tripleLookup rho x
                           | Just w <- lookup rho x
                           , Just y <- lookup rho w
                           , Just z <- lookup rho y
                           = z
                        tripleLookup _ x = handleFailure x
        \end{verbatim}
        \end{center}    
    \caption{Pattern guards swoop in to restore the Nice~Properties, and all is
    right again.} 
    \Description{An implementation of nameOf using pattern guards.}
    \label{fig:guardtriplelookup}
    \end{figure}


    \rab{Or pattern example}
    Or patterns do \dots which eliminate the need to \dots

     \begin{figure}
            \begin{center}
                \begin{verbatim}
        let game_of_token token = match token with 
            | BattlePass f      -> ("Fortnite", f)
            | ChugJug f         -> ("Fortnite", f)
            | TomatoTown f      -> ("Fortnite", f)
            | HuntersMark f     -> ("Bloodborne", 2 * f)
            | SawCleaver f      -> ("Bloodborne", 2 * f)
            | MoghLordOfBlood f -> ("Elden Ring", 3 * f)
            | PreatorRykard f   -> ("Elden Ring", 3 * f)
            | _                 -> ("Irrelevant", 0)
                \end{verbatim}
            \end{center}    

        \caption{\tt{game\_of\_token}, with redundant right-hand sides,
        should raise a red flag.} 
        \Description{An implementation of game_of_token using bare patterns.}
        \label{fig:baregot}
        \end{figure}


          \begin{figure}
    \begin{center}
    \begin{verbatim}
let game_of_token token = match token with 
    | BattlePass f | ChugJug f | TomatoTown f  -> ("Fortnite", f)
    | HuntersMark f | SawCleaver f             -> ("Bloodborne", 2 * f)
    | MoghLordOfBlood f | PreatorRykard f      -> ("Elden Ring", 3 * f)
    | _                                        -> ("Irrelevant", 0)
    \end{verbatim}
    \end{center}    
    \caption{Or-patterns condense~\tt{game\_of\_token}
    significantly, and it is easier to read line-by-line.} 
    \Description{An
    implementation of game_of_token using or-patterns.}
    \label{fig:orgot}
    \end{figure}

    \newpage
    Individually, each of these three extensions grants a powerful abstraction
    that reduces code duplication. It follows that they should retain their
    powers of brevity when operating in unison, and indeed this is the case:
    with these three extensions co-operating, the previous example reduces like
    this \dots

    Given the expressive power of this combination and each extension's
    individual popularity, it would seem obvious that they are all supported by
    myriad functional programming languages. But that's not true: among Haskell,
    OCaml, Scala, Erlang/Elixir, Rust, F\#, Rocq, and Agda~\citep{haskell,
    ocaml, scala, erlang, elixir, rust, fsharp, agda}, only Haskell has all
    three of side conditions, or patterns, and pattern guards- and only as of
    2023! In fact, when the original idea for this work was conceived in 2023, a
    primary motivator for suggesting an alternative to pattern matching was that
    there was \textit{no} language with all three of them. 

    - No strict language with all three + significant implementation effort in
    Haskell to get them all working -> evidence enough that having all of them
    is a non-trivial problem. 
    
    We~find the extension story somewhat unsatisfying. At the very least,
    we~want to be able to use pattern matching, with the extensions we~want, in
    a single language. Or, we~want an alternative that gives me the expressive
    power of pattern matching with these extensions. 

    \section{Equations}

    An intriguing alternative to pattern matching exists in~\it{equations}, from
    the Verse Calculus (\VC), a core calculus for the functional logic
    programming language~\it{Verse}~\citep{antoy2010functional,
    hanus2013functional, verse}. (For the remainder of this work, we~use “\VC”
    and “Verse” interchangeably.)

    Verse's equations and choice mechanism allow for much of the same expressive
    power as pattern matching with extensions. Rather than go into too much
    detail (the curious may examine the Verse paper, which is well worth a
    read), we support this claim with a familiar example: 

    \rab{TODO : insert existing examples}

    The brevity of the code in the figures is is promising for~\VC. If it rivals
    pattern matching with popular extensions in expressiveness, and~\VC
    does everything using only equations and choice, it seems like the language
    is an intriguing option for writing code! 


    But the VC semantics has a cost model that is hard to predict. Why? In~\VC,
    names (logical variables) are~\it{values}, and they can just as easily be
    the result of evaluating an expression as an integer or tuple. To bind these
    names,~\VC, like other functional logic languages, relies
    on~\it{unification} of its logical variables and \it{search} at runtime to
    meet a set of program constraints~\citep{antoy2010functional,
    hanus2013functional}. Combining unifying logical variables with search
    requires backtracking, which can lead to exponential runtime
    cost~\citep{hanus2013functional, wadler1985replace, clark1982introduction}. 

    % Todo: An example would be nice. Backtracking/evaluating backwards, then
    % done. 

    Pattern matching, by contrast, can be compiled to a~\it{decision tree}, a
    data structure that enforces linear runtime performance by guaranteeing no
    part of the scrutinee will be examined more than once~\citep{maranget}. A
    decision tree does not backtrack: once a program makes a decision based on
    the form of a value, it does not re-test it later with new information. 
    

    Figures: 

    \begin{figure}[ht] 
        \begin{minipage}[h]{0.54\linewidth}
          \verselst
          \begin{lstlisting}[numbers=none, basicstyle=\tiny, xleftmargin=.2em,
                            showstringspaces=false,
                            frame=single]
$\exists$ exclaimTall. exclaimTall = $\lambda$ sh. 
  one { 
      $\exists$ s. sh = $\langle$Square s$\rangle$; 
      s > 100.0; "Wow! That's a sizeable square!"
    | $\exists$ w h. sh = $\langle$Triangle, w, h$\rangle$; 
      h > 100.0; "Goodness! Towering triangle!"
    | $\exists$ b1 b2 h. sh = $\langle$Trapezoid, b1, b2, h$\rangle$;
      h > 100.0; "Zounds! Tremendous trapezoid!"
    | "Your shape is small." }
            \end{lstlisting}
                \Description{exclaimTall}
        \subcaption{\tt{exclaimTall} in~\VC }
            \label{fig:verseexclaimtall} 
        %   \vspace{4ex}
        \end{minipage}%%
        \begin{minipage}[h]{0.5\linewidth}
          \verselst
          \begin{lstlisting}[numbers=none, basicstyle=\tiny, xleftmargin=2em,
                        frame=single]
$\exists$ tripleLookup. tripleLookup = $\lambda$ rho x. 
    one { $\exists$ w. lookup rho x = $\langle$Just w$\rangle$; 
          $\exists$ y. lookup rho w = $\langle$Just y$\rangle$; 
          $\exists$ z. lookup rho y = $\langle$Just z$\rangle$;
          z
        | handleFailure x }
          \end{lstlisting}
                \Description{exclaimTall}
            \subcaption{\tt{tripleLookup} in~\VC }
            \label{fig:versetriplelookup} 
        \vspace{4ex}
        \end{minipage} 
        \begin{minipage}[h]{\linewidth}
          \verselst
          \begin{lstlisting}[numbers=none, basicstyle=\tiny, xleftmargin=9em, 
                            frame=single, showstringspaces=false]
$\exists$ game_of_token. game_of_token = $\lambda$ token. 
  $\exists$ f. one {
         token = one { $\langle$BattlePass, f$\rangle$ | $\langle$ChugJug, f$\rangle$ | $\langle$TomatoTown, f$\rangle$}; 
           $\langle$"Fortnite", f$\rangle$
       | token = one { $\langle$HuntersMark, f$\rangle$ | $\langle$SawCleaver, f$\rangle$}; 
           $\langle$"Bloodborne", 2 * f$\rangle$
       | token = one { $\langle$MoghLordOfBlood, f$\rangle$ | $\langle$PreatorRykard, f$\rangle$}; 
           $\langle$"Elden Ring", 3 * f$\rangle$
       |   $\langle$"Irrelevant", 0$\rangle$ }
          \end{lstlisting}
                \Description{game_of_token in~\VC}
        \subcaption{\tt{game\_of\_token} in~\VC }
            \label{fig:versegot} 
        \vspace{4ex}
        \end{minipage}%% 
    \caption{Code for the~\ref{pmandextensions} functions with equations looks
    similar, and it doesn't need extensions.}
    \label{fig:verseextfuncs}
      \end{figure}
        


      \rab{todo: name this better}
    \section{Two Languages, Two Solutions}
    We have created and implemented a semantics for three core languages that
    aim to bridge the gap between pattern matching, equations, and decision
    trees. They are \PPlus, which has pattern matching with all the above
    extensions, \VMinus, which has equations without backtracking or multiple
    results, and \D, which has decision trees. \PPlus\ and \VMinus serve to
    explore the design space of pattern matching and equations, while \D\ serves
    as the target of translation to reason about efficiency. Each language is an
    extension of the untyped lambda calculus with value constructors. 
    
    Like in \VC, every lambda calculus term is valid in our languages and has
    the same semantics. Also like the lambda calculus and \VC, all three
    languages are \it{strict}, meaning expressions are always evaluated when
    they are bound to variables, and they are all untyped. 
    
    % In the subsections below, we discuss each language in more detail. Section 7
    % discusses futher how \PPlus\ and \VMinus\ relate to modern languages with
    % pattern matching and \VC\ respectively. 

\subsection{Introducing \PPlus}
\label{pplus}

    \PPlus\ packages common extensions to pattern matching with a standard core.
    In addition to bare pattern matching— names and applications of value
    constructors— the language includes pattern guards, or-patterns, and side
    conditions. You may point out that pattern guards subsume side conditions
    (yes, they do!); we include side conditions and separate them from guards
    for purpose of study. 
    
    % Furthermore, \PPlus\ introduces a new form of pattern:
    % \pcommap. This allows a pattern in \PPlus\ be a \it{sequence} of
    % sub-patterns, allowing a programmer to stuff as many patterns as they want
    % in the space of a single oneWe discuss the implications of this in
    % Section~\ref{ppweird}.

\subsection{Formal Semantics of \PPlus}

% \begin{figure}
% \begin{center}
% \begin{bnf}
% $P$ : \textsf{Programs} ::=
% $\bracketed{d}$ : definition
% ;;
% $d$ : \textsf{Definitions} ::=
% | $\tt{val}\; x\; \tt{=}\; \expr$ : bind name to expression
% ;;
% $\expr$ : Expressions ::= 
% | $v$ : literal values 
% | $x, y, z$ : names
% | $K\bracketed{\expr}$ : value constructor application 
% | $\ttbackslash x.\; \expr$ : lambda declaration  
% | $\expr[1]\; \expr[2]$ : function application 
% | $\tt{case}\; \expr\; \bracketed{p\; \ttrightarrow\; \expr}$ : case expression 
% | \ttbraced{$\expr$}
% ;;
% $p$ : \textsf{Patterns} ::= $p_{1}\pbar p_{2}$ : or-pattern
% | $p \tt{,} p'$ : pattern guard 
% | $p\; \tt{<-}\; \expr$ : pattern from explicit expression  
% | $x$ : name 
% | $\tt{\_}$ : wildcard 
% | $K\; \bracketed{p}$ : value constructor application 
% | $\tt{when}\; \expr$
% | \ttbraced{$p$}
% ;;
% $\v$ : Values ::= $K\bracketed{\v}$ : value constructor application 
% | $\ttbackslash x.\; \expr$ : lambda value 
% ;;
% $K$ : \textsf{Value Constructors} ::=
% | $\tt{true}\; \vert\; \tt{false}$ : booleans
% | $\tt{A-Z}x$ : name beginning with capital letter
% % | $[\tt{-}\vert\tt{+}](\tt{0}-\tt{9})+$ : signed integer literal 
% \end{bnf}
% \end{center}
% \Description{The concrete syntax of PPlus.}
% \caption{\PPlus: Concrete syntax}
% \label{fig:ppsyntax}
% \end{figure}

%   % \subsection{Forms of Judgement for {\PPlus}:}
% \begin{tabular}{ll}
% \toprule
%     \multicolumn2{l}{\emph{Metavariables}} \\
% \midrule
%     $\expr$& expression \\
%     $\v,\; \v'$& value \\
%     $K$& value constructor \\ 
%     $p$& pattern \\ 
%     $x, y$& names \\ 
%     $g$& pattern guard \\ 
%     $gs$& list of pattern guards \\ 
%     $\dagger$& pattern match failure \\ 
%     $s$& a solution, either $\v$ or $\dagger$ \\ 
%     \Rho& environment: $name \rightarrow \mathcal{\v}$ \\
%     \Rho\:+\:\Rhoprime& extension \\
%     $\Rho \uplus \Rho'$& disjoint union \\
%     $\{ x \mapsto y \} $& environment with name $x$ mapping to $y$ \\
% \bottomrule
% \end{tabular}    


% \begin{align*}
%     &\pmatch \quad   \rm{ Pattern $p$ matches value $\v$, 
%                           producing bindings $\rho'$;} \\
%     &\pfail  \quad\; \rm{ Pattern $p$ does not match value $\v$.} 
% \end{align*}

% If a pattern is bound to an expression in the form $\pmatch[pat = {p \leftarrow
% \expr}, newenv = \solution]$, it will match if the expression $e$ evaluates to
% value $v'$ and $p$ matches with $v'$. If a pattern is standalone, as in any
% other form, it will match on $v$, the \it{original} scrutinee of the case
% expression. For example, in $\pmatch[pat = {p_{1}, p_{2} \leftarrow \expr},
% newenv = \solution]$, $p_{1}$ matches with $v$, and when $\ppeval[result=\v']$,
% $p_{2}$ matches with $\v'$. 

% Pattern matching is defined inductively: 
% \begin{itemize}
%     \item a name $x$ matches any value $\v$, and produces the environment 
%     $\bracketed{x \mapsto \v}$. 
%     \item a value constructor $K$ applied to atoms  matches 
%     a value $\v$ if $\v$ is an application of $K$ to the same number of values,
%     each of which matches the corresponding atom. Its match produces 
%     the disjoint union of matching all internal atoms to all internal values. 
%     \item a pattern \whenexpr\ matches when $\expr$ evaluates to a value other than 
%     $\mathit{false}$, and produces the empty environment. 
%     \item a pattern \parrowe\ matches when $\expr$ evaluates to 
%           value $\v$, and $p$ matches $\v$. 
%     \item a pattern \pcommap\ matches if both $p$ and $p_{2}$ match.
%     \item a pattern \porp\ matches if either $p_{1}$ or $p_{2}$
%     matches. 
% \end{itemize}

% When a pattern is of the form $K, p_{1}, \dots, p_{n}$, each sub-pattern $p_{i}$
% may introducing new variables during pattern matching. Bindings for all these
% variables must be combined in a single environment. \it{Disjoint union} is an
% operation on two environments. Its definition borrowed, paraphrased, from Ramsey\cite{bpc}" 

% The disjoint union of environments $\rho_{1}$ and $\rho_{2}$ is written
% $\rho_{1} \uplus \rho_{12}$, and it is defined if and only if $\dom \rho_{1}
% \cap \dom \rho_{2} = \emptyset$:

% \begin{align*}
%   \dom(\rho_{1} \uplus \rho_{2}) &= \dom \rho_{1} \uplus dom \rho_{2}, \\
%     (\rho_{1} \uplus \rho_{2})(x) &= 
%   \begin{cases}
%     \rho_{1}(x), & \text{if } x \in \dom  \rho_{1} \\
%     \rho_{2}(x), & \text{if } x \in \dom \rho_{2} \\
% \end{cases}
% \end{align*}

% Disjoint union across multiple results, when any result is failure, still
% represents failure: 

% \begin{gather*}
%   \dagger \uplus \rho = \dagger \quad
%   \rho \uplus \dagger = \dagger \quad
%   \dagger \uplus\, \dagger = \dagger
% \end{gather*}

% Finally, when a pattern in a branch in a case expression fully matches, its
% corresponding right-hand side is evaluated with top-level $\rho$ extended with
% the $\rho'$ produced by the pattern match. Environment extension is notated
% $\rho + \rho'$. 

% \subsection{Expressions in \PPlus}

%     An expression in {\PPlus} evaluates to produce a single value. 
%     \begin{align*}
%         &\prun
%     \end{align*}
    
%     \ppsemantics 
      
%       \bigskip 

    
  Figure~\ref{fig:ppfuncs} provides an example of how a programmer might utilize
  \PPlus\ to solve the previous problems: 

    \begin{figure}[ht] 
      \begin{minipage}[h]{0.54\linewidth}
        \pplst 
        \begin{lstlisting}[numbers=none, basicstyle=\tiny, xleftmargin=.2em,
          showstringspaces=false,
          frame=single]
val exclaimTall = \sh.
  case sh of Square s, when (> s) 100 -> 
            Wow! That's A Tall Square!  
    | Triangle w h, when (> s) 100 ->
            Goodness! Towering Triangle!
    | Trapezoid b1 b2 h, when (> s) 100 -> 
            Zounds! Tremendous Trapezoid!
    | _ ->  Your Shape Is Small
  \end{lstlisting}
      \Description{exclaimTall in PPlus}
      \subcaption{\tt{exclaimTall} in \PPlus}
          \label{fig:ppexclaimtall} 
      %   \vspace{4ex}
      \end{minipage}%%
      \begin{minipage}[h]{0.5\linewidth}
        \pplst 
        \begin{lstlisting}[numbers=none, basicstyle=\tiny, xleftmargin=2em,
                      frame=single]
  val tripleLookup = \ rho. \x.
    case x of 
      Some w <- (lookup rho) x
    , Some y <- (lookup rho) w
    , Some z <- (lookup rho) y -> z
    | _ -> handleFailure x

   \end{lstlisting}
        \Description{tripleLookup in PPlus}
        \subcaption{\tt{tripleLookup} in \PPlus}
            \label{fig:pptriplelookup} 
        \vspace{4ex}
      \end{minipage} 
      \begin{minipage}[h]{\linewidth}
        \pplst 
        \begin{lstlisting}[numbers=none, basicstyle=\tiny, xleftmargin=9em,
          showstringspaces=false,
          frame=single]
val game_of_token = \t. 
  case t of  
    BattlePass f  | (ChugJug f | TomatoTown f) -> P (Fortnite f)
  | HuntersMark f | SawCleaver f               -> P (Bloodborne ((* 2) f))
  | MoghLordOfBlood f | PreatorRykard f        -> P (EldenRing ((* 3) f))
  | _                                          -> P (Irrelevant 0)
\end{lstlisting}
      \Description{game_of_token in PPlus}
      \subcaption{\tt{game\_of\_token} in \PPlus}
          \label{fig:ppgot}
      \vspace{4ex}
      \end{minipage}%% 
      \caption{The functions in \PPlus\ have the same desirable implementation
      as before.}
  \label{fig:ppfuncs}
    \end{figure}        
    
    % We've had to make some syntactic concessions: \PPlus doesn't have strings or
    % tuples, so we've used masquerading value constructors, and the parser has no
    % left recursion or notion of infix operators, so we've had to parenthesize a
    % few expressions. Extending \PPlus\ with strings, left recursion, true
    % tuples, and math operators is left as an exercise to the reader. 
    
% \subsection{Addressing how \PPlus\ handles unusual pattern combinations}
% \label{ppweird}
%     Given its minimal restrictions on what kind of pattern can appear where,
%     \PPlus\ admits of strange looking patterns: consider \tt{Cons (when true)
%     zs}. But these should not be alarming, because such syntactic forms reduce
%     to normal forms by (direct) application of algebraic laws: 

%     \begin{align}
%       K (\whenexpr) \;p2\; \dots &=== K \;\_ \;p2\; \dots, \whenexpr \\
%       K (\whenexpr, \;p2\;) \;p3\; \dots  &=== K \;p2\; \;p3\; \dots, \whenexpr \\
%       K (p1\;, \whenexpr) \;p3\; \dots  &=== K \;p1\; \;p3\; \dots, \whenexpr \\
%       K (\whenexpr \pbar p2\;) \;p3\; \dots &=== (K \;\_ \;p3\; \dots, \whenexpr) \pbar (K \;p2\; \;p3\; \dots) \\
%       K (p1 \pbar \whenexpr) \;p3\; \dots &=== (K \;p2\; \;p3\; \dots) \pbar (K \;\_ \;p3\; \dots, \whenexpr)  \\
%       \whenexpr \leftarrow e &=== \;\_ <- e, \whenexpr
%     \end{align}
    
    
%     Repeatedly applying these laws until the program reaches a fixed point
%     normalizes placements of \it{when}. Laws (2) and (3) work because \PPlus\
%     has no side effects and the laws assume all names are unique (the compiler
%     takes care of this), so changing the order in which patterns match has no
%     effect on semantics.         


\subsection{Introducing \VMinus\ }
\label{vminus}

        To fuel the pursuit of smarter decision-making, we now draw inspiration
        from \VC. We begin by trimming away the aforementioned culprits that
        tends to lead to unpredictable or costly outcomes during runtime:
        backtracking and multiple results. Removing these language traits strips
        much of the identity of \VC; however, we do so in order to study only
        \VC's \it{equations} while staying grounded in an otherwise-typical
        programming context of single results and no backtracking at runtime.
        Both backtracking and multiple results often manifest within \VC's
        choice operator, which tempt its removal altogether. However, we are drawn
        to harnessing the expressive potential of this operator, particularly
        when paired with \VC's 'one' keyword. When employed with choice as a
        condition, 'one' elegantly signifies 'proceed if any branch of the
        choice succeeds. 
        
        To this end, we imagine a new language, \VMinus, in which choice is
        permitted with several modifications:

        \begin{enumerate}
        \item Choice may only appear as a condition or 'guard', not as a result
        or the right-hand side of a binding.
        \item If any branch of the choice succeeds, the choice succeeds,
        producing any bindings found in that branch. The program examines the
        branches in a left-to-right order.
        \item The existential $\exists$ may not appear under choice.
        \end{enumerate}

        We introduce one more crucial modification to the \VC\ runtime: a name in
        \VMinus is an \it{expression} rather than a \it{value}. This alteration,
        coupled with our adjustments to choice, eradicates backtracking. Our
        rationale behind this is straightforward: if an expression returns a
        name, and subsequently, the program imposes a new constraint on that
        name, it may necessitate the reevaluation of the earlier expression—- a
        scenario we strive to avoid. 
        
        % An example in \VC\ illustrates this precise
        % scenario: 

        % FIRST PAPER EXAMPLE 

        % IMPORTANTLY, \VMinus REALLY ISN'T A FUNCTIONAL LOGIC PROGRAMMING LANGUAGE ANYMORE. 
        
        %  Doesn't backtrack
        
        %  No multiple results 
        
        %  Doesn't evaluate functions backwards, have top-level patterns like verse, the list goes on 

        % If you're well-versed in functional logic programming, may perceive that
        % imposing such restrictions on choice and names effectively strips away
        % much of Verse's essence as a functional logic programming language. With
        % these constraints enforced, there can be no backtracking, multiple
        % results, backward function evaluation, or top-level patterns, among
        % other classic functional logic programming features. But do not fear:
        % our intent is not to recklessly strip Verse of its meticulously crafted
        % core tenets. Instead, our aim is to extract a select few—namely, its
        % equations, existentials, and nondeterministic evaluation order—and
        % juxtapose them with pattern matching.    

        Indeed, having stripped out the functional logic programming elements of
        \VC, only the decision-making bits are left over. To wrap these, we add
        a classic decision-making construct: guarded commands\cite{dijkstra} The
        result is \VMinus. \VMinus's concrete syntax is defined in
        Figure~\ref{fig:vmsyntax}. 

        % \bf{Choice forces "name knowing." Names must be known in all branches 
        % if bound to unknown variable or bound to known variable if any name 
        % is unknown}.
        
        % With these modifications in place, alongside a few additional
        % adjustments to the placement of existentials and the timing of choice
        % evaluation (footnote: we discuss these later), we unveil a redefined
        % decision-making core. To complete this transformation, we incorporate a
        % venerable decision-making construct: Guarded commands [cite: Dijkstra].
        % The result is \VMinus. 
        % \begin{figure}[H]
        %   \begin{center}
        %       \begin{bnf}
        %       $P$ : \textsf{Programs} ::=
        %       $\bracketed{d}$ : definition
        %       ;;
        %       $d$ : \textsf{Definitions} ::=
        %       | $\tt{val}\; x\; \tt{=}\; \expr$ : bind name to expression
        %       ;;
        %       $\expr$ : \textsf{Expressions} ::=
        %       | $v$ : literal values 
        %       | $x, y, z$ : names
        %       | $\tt{if}\; [\; G\; \bracketed{\tt{[]}\; G}\;]\; \tt{fi}$ : if-fi 
        %       | $K \bracketed{\expr}$ : value constructor application 
        %       | $\expr[1]\; \expr[2]$ : function application 
        %       | $\ttbackslash x\tt{.}\; \expr$ : lambda declaration 
        %       ;;
        %       $G$ : \textsf{Guarded Expressions} ::=  
        %       $[\tt{E }{\bracketed{x}}{\tt{.}}] \bracketed{g} \;\ttrightarrow\; \expr$ : names, guards, and body
        %       ;;
        %       $g$ : \textsf{Guards} ::=  
        %       | $\expr$ : intermediate expression 
        %       | $x \;\tt{=}\; \expr$ : equation 
        %       | $ g \bracketed{\tt{;} g} \; \pbar \; g \bracketed{\tt{;} g}$ : choice 
        %       ;;
        %       $\v$ : Values ::= $K\bracketed{\v}$ : value constructor application 
        %       | $\ttbackslash x\tt{.}\; \expr$ : lambda value
        %       \end{bnf}
        %   \end{center}
        %   \Description{The concrete syntax of VMinus.}
        %       \caption{\VMinus: Concrete syntax}
        %       \label{fig:vmsyntax}
        %       \end{figure}      

    \bf{Programming in \VMinus}
        
    Even with multiple modifications, \VMinus\ still allows for many of the same
    pleasing computations as full Verse. Programmers in \VMinus can\dots
        \begin{enumerate}
            \item Introduce a set of equations, to be solved in a
            nondeterministic order 
            \item Guard expressions with those equations 
            \item Flexibly express “proceed when any of these operations
            succeeds” with the new semantics of choice.  
        \end{enumerate}

    Figure~\ref{fig:vminusfuncs} provides an example of how a programmer can
    utilize \VMinus\ to solve the previous problems: 
    
    
    \begin{figure}[ht] 
      \begin{minipage}[h]{0.54\linewidth}
        \vmlst 
        \begin{lstlisting}[numbers=none, basicstyle=\tiny, xleftmargin=.2em,
          showstringspaces=false,
          frame=single]
val exclaimTall = \sh.
  if sh = Square s; (> s) 100 -> 
            Wow! That's A Tall Square!  
    [] E w h. sh = Triangle w h; (> s) 100 ->
            Goodness! Towering Triangle!
    [] E b1 b2 h. sh = Trapezoid b1 b2 h; 
     (> s) 100 -> Zounds! Tremendous Trapezoid!
    [] ->  Your Shape Is Small
  fi 
  \end{lstlisting}
      \Description{exclaimTall in VMinus}
      \subcaption{\tt{exclaimTall} in \VMinus}
          \label{fig:vmexclaimtall} 
      %   \vspace{4ex}
      \end{minipage}%%
      \begin{minipage}[h]{0.5\linewidth}
        \vmlst 
        \begin{lstlisting}[numbers=none, basicstyle=\tiny, xleftmargin=2em,
                      frame=single]
  val tripleLookup = \ rho. \x.
    if E x w y z v1 v2 v3. 
      v1 = (lookup rho) x; v1 = Some w; 
      v2 = (lookup rho) w; v2 = Some y; 
      v3 = (lookup rho) y; v3 = Some z; -> 
            z 
      [] -> handleFailure x
    fi 
   \end{lstlisting}
        \Description{tripleLookup in VMinus}
        \subcaption{\tt{tripleLookup} in \VMinus}
            \label{fig:vmtriplelookup} 
        \vspace{4ex}
      \end{minipage} 
      \begin{minipage}[h]{\linewidth}
        \vmlst 
        \begin{lstlisting}[numbers=none, basicstyle=\tiny, xleftmargin=9em,
          showstringspaces=false,
          frame=single]
val game_of_token = \t. 
  if E f. (t = BattlePass f  | t = ChugJug f | t = TomatoTown f) -> 
        P (Fortnite f)
  [] E f. (t = HuntersMark f | t = SawCleaver f)                 -> 
        P (Bloodborne ((* 2) f))
  [] E f. (t = MoghLordOfBlood f | t = PreatorRykard f)          -> 
        P (EldenRing ((* 3) f))
  [] -> P (Irrelevant 0)
  fi 
\end{lstlisting}
      \Description{game_of_token in VMinus}
      \subcaption{\tt{game\_of\_token} in \VMinus}
          \label{fig:vmgot}
      \vspace{4ex}
      \end{minipage}%% 
      \caption{The functions in \VMinus\ have a desirably concise
      implementation, as before.}
  \label{fig:vminusfuncs}
    \end{figure}   

    \VMinus\ looks satisfyingly similar to both \PPlus\ and \VC. 

    \rab{include section?}
    \subsection{\VMinus\ and \PPlus, side by side}
    % We compare \VMinus\ with \PPlus\ as an exercise in comparing equations with
    % pattern matching. They certainly look similar, which hints that \VMinus\
    % might be as expressive as pattern matching with the three extensions. 
    % Proving this claim is the topic of Section~\ref{pplustovminus}. 

    % When might a programmer prefer \VMinus\ over \PPlus, or vice versa? After
    % programming in both, we have come up with three empirical observations: 

    % \begin{enumerate}
    %   \item \it{The scrutinee:} When there is no obvious single scrutinee, \VMinus is more succinct.
    %   When there is a scrutinee, \PPlus is more succinct. 
    %   \item \it{Binding and decision-making:} Binding and decision-making are
    %   joined in a single construct in \VMinus: \tt{=}. \PPlus\ needs different
    %   kinds of syntax, like \tt{<-}, to express different kinds of binding. And
    %   a programmer wanting for \tt{let} in \PPlus\ will never feel this need in
    %   \VMinus: \tt{=} subsumes that, too. 
    %   \item \it{Names:} In \VMinus, names are explicitly introduced. In
    %   practice, this helps prevent a common mistake in pattern matching in which
    %   a programmer attempts to match a value $v$ on an in-scope name $x$,
    %   expecting the match to succeed iff $x$ evaluates to $v$ at runtime, only
    %   to see the match always succeed. 
    % \end{enumerate}     

    % \rab{How do you feel about examples for each of these points here? This
    % section is very long, but they help drive these points home.}

% \subsection{Concrete Syntax of \VMinus}

% \bigskip

% \begin{center}
%     \begin{bnf}
%     $P$ : \textsf{Programs} ::=
%     $\bracketed{d}$ : definition
%     ;;
%     $d$ : \textsf{Definitions} ::=
%     | $\tt{val}\; x\; \tt{=}\; \expr$ : bind name to expression
%     ;;
%     $\expr$ : \textsf{Expressions} ::=
%     | $v$ : literal values 
%     | $x, y, z$ : names
%     | $\tt{if}\; \tt{[}\; G\; \bracketed{\square\; G}\; \tt{]}\; \tt{fi}$ : if-fi 
%     | $K \bracketed{\expr}$ : value constructor application 
%     | $\expr[1]\; \expr[2]$ : function application 
%     | $\lambda x.\; \expr$ : lambda declaration 
%     ;;
%     $G$ : \textsf{Guarded Expressions} ::=  
%     $[\vexists{\bracketed{x}}] \bracketed{g} \boldsymbol{\rightarrow}\expr$ : names, guards, and body
%     ;;
%     $g$ : \textsf{Guards} ::=  
%     | $\expr$ : intermediate expression 
%     | $x = \expr$ : equation 
%     | $ g \bracketed{; g} \; \choice \; g \bracketed{; g}$ : choice 
%     ;;
%     $\v$ : Values ::= $K\bracketed{\v}$ : value constructor application 
%     | $\lambda x.\; \expr$ : lambda value
%     ;;
%     \end{bnf}
% \end{center}

% A \it{name} is any token that is not an integer literal, 
% does not contain whitespace, a bracket, or parenthesis, 
% and is not a value constructor name or a reserved word.

\subsection{Formal Semantics of ${\VMinus}$:}
% \begin{tabular}{ll}
% \toprule
%     \multicolumn2{l}{\emph{Metavariables}} \\
% \midrule
%     $\expr$& An expression \\ 
%     $\v,\; \v'$& value \\
%     $\fail$& expression failure \\
%     $\result$& $\v \vert \fail$ : expressions produce \it{results}: values or
%     failure. \\
%     \Rho& environment: $name \rightarrow {\mathcal{V}}_{\bot}$ \\
%     $\rho\{ x \mapsto y \} $& environment extended with name $x$ mapping to $y$ \\
%     $g$& A guard \\
%     $eq$& equation \\ 
%     % $\tempstuck$& a temporarily-stuck equation \\
%     $\dagger$& when solving guards is rejected \\
%     $\result$& $\rhohat \mid \dagger$ : guards produce \it{solutions}: a
%     refined environment \rhohat\; or rejection\\
%     $\mathcal{T}$& Context of all temporarily stuck guards (a sequence) \\ 
%     $G$& A guarded expression \\
%     % \uppsidown& Inability to compile to a decision tree; a compile time error \\
% \bottomrule
% \end{tabular}    

\bigskip
    
    \subsection{Expressions}
    
    \newcommand\GNoTree{\vmrung \rightsquigarrow \uppsidown} 
    
    An expression in core Verse evaluates to produce possibly-empty sequence of
    values. In {\VMinus}, expressions evaluate to a single value. 
    % We will explore a Verse variant in which expressions may evaluate to
    % multiple values. 

    An expression evaluates to produce a \bf{result}. A result is either a
    value $\v$ or \fail. 
    
    \[\result\; \rm{::=}\; \rhohat\; \vert \; \fail \]
    
    \showvjudgement{Eval}{\vmeval}
    \showvjudgement{Process-Guards}{\vmgs}
\subsection{Guards}

In {\VMinus}, we solve guards similarly to how the authors of the original Verse
paper solve equations: we pick one, attempt to solve for it, and move on. 

In our semantics, we do this with a \it{list} of guards: we pick a guard,
attempt to solve it to refine our environment or fail, and repeat. We only pick
a guard out of the context \context\ that we know we can solve. If we can't pick
a guard and there are guards left over, the semantics gets stuck.

Given an environment $\rho$ from names to values {\v}s plus $\bot$, solving a
guard \g\ will either refine the environment ($\Rhoprime$) or be \bf{rejected}.
We use the metavariable $\dagger$ to represent rejection, and a guard producing
$\dagger$ will cause the top-level list of guards to also produce $\dagger$. 

    \showvjudgement{Guard-Refine}{\grefine[result=\rho']}
    \showvjudgement{Guard-reject}\gfail

%     \subsection{Refinement ordering on environments}

%     \begin{align*}
%     \rho \subseteq \Rhoprime \text{ when }&\dom\rho  \subseteq \dom \Rhoprime\\
%     \text{ and } &\forall x \in \dom \rho: \rho(x) \subseteq \Rhoprime(x)
%     \end{align*}
    
%     \medskip
        
%     Success only refines the environment; that~is, when ${\vtuple{\rho, \expr}}
%     \rightarrowtail{} {\Rhoprime}$, we~expect $\rho \subseteq \Rhoprime$.
    
%     \vfilbreak

\rab{I put vminus semantics here, but they're quite long. How much to show?}
% \vmsemantics

% \subsection{Introducing \D}
% \label{d}

% While \PPlus\ and \VMinus exist for writing programs, \D\ exists as the target
% of translation. This is because the decision-making construct in \D, the
% \it{decision tree}, is impractical for programmers but practical for theorists
% who wish to prove properties of efficiency. Specifically, a decision tree is an
% automaton that implements pattern matching with the invariant that no part of
% the scrutinee of a \it{case} expression is examined more than once at runtime.
% Through this invariant, a decision tree guarantees linear runtime performance of
% pattern matching. This desirable property of decision trees is paired with the
% disadvantage of potential code size explosion. 

% Decision trees are classically used as an intermediate representation for
% compiling \it{case} expressions. In this work, we will use them as a target for
% compiling \it{if-fi} in \VMinus. We do so to prove that equations, at least given
% their restrictions in \VMinus, can be compiled to a decision tree. 

% \D\ is a simple generalization of the trees found in Maranget\cite{maranget}We
% discuss the specifics of \D's decision trees in the following section. 

\subsection{\VMinus compiles to a decision tree in the ~\D language}
\label{d}

Like pattern matching, \VMinus~compiles to a decision tree whose height is 
linear with respect to code size. To formalize decision trees, we use 
\D, a language which generalizes the trees found in~\citet{maranget}. 

The spec of decision trees and match compilation algorithm is presented below. 

% While \PPlus\ and \VMinus exist for writing programs,~\D exists as the target of
% translation and provides a means by which to demonstrate~\VMinus's efficient
% cost model. This is because the decision-making construct itself in~\D, the
% \it{decision tree}, has an efficient cost model. A decision tree can implement
% either pattern matching or~\iffibf while guaranteeing never to repeat a test.
% The worst-case cost of evaluating a decision tree is linear in its depth, which
% itself linear in the size of the code. This desirable property of decision trees
% is half of a space-time tradeoff: when a decision tree is produced by compiling
% a~\it{case} expression, there are pathological cases in which the total size of
% the tree is exponential in the size of the source code (from~\it{case}). Run
% time remains linear, but code size may not be. 

% Although decision trees are classically used as an intermediate representation
% for compiling~\it{case} expressions, in this work, we~use them as a target for
% compiling~\iffibf in~\VMinus. In particular, we~show that equations in~\VMinus
% can be compiled to a decision tree. 

% \D is a generalization of the trees found in~\citet{maranget}. 

% \subsection{\D is a generalization of Maranget's trees} 

% \D's syntax is given in Figure~\ref{fig:dsyntax}. Decision trees in~\D are
% engineered to look like Maranget's trees. 
  
      
%   The heart of a decision tree is the~\it{test} form: it takes a value, examines
%   it, and chooses a branch based on its form (Maranget calls the operation
%   \textsc{Switch}). 

%   Let's look at an example from~\citet{maranget} which shows the structure of a
%   simple pattern-matching function and its corresponding decision treeWe show
%   Maranget's example for two reasons: First, the example serves to bolster your
%   understanding of decision trees with a classic, well-established model.
%   Second, since we currently have no visualization generator for~\D (\D~can
%   currently only be visualized as plain text), we use Maranget's example to have
%   a reasonable visual representation of decision trees. The example beings with
%   the function,~\tt{merge}, which merges two lists: 

%   \begin{figure}[H]
%       \includegraphics[scale=0.7]{../images/merge.png}
%       \Description{Maranget's merge implementation}
%       \caption{The skeleton of Maranget's~\tt{merge}}
%   \end{figure}

  % He then shows an intermediate representation in his compilation algorithm,
  % the occurrence vector and clause matrix: 

  %~\begin{figure}[H]
  %     \begin{gather*}
  %         \vec{o} = (\tt{xs ys})~\hspace{3em}
  %         P~\rightarrow A = 
  %         \begin{pmatrix}
  %            ~\tt{[]}     &~\tt{\_}     &~\rightarrow~\tt{xs}~\\
  %            ~\tt{\_}     &~\tt{[]}     &~\rightarrow~\tt{ys}~\\
  %            ~\tt{\_::\_} &~\tt{\_::\_} &~\rightarrow~\tt{...} 
  %         \end{pmatrix}
  %     \end{gather*}


  %     \Description{Occurrence vector of names, and clause matrix of matches}
  %     \caption{The occurrence vector and clause matrix are intermediate
  %     representations in Maranget's compilationWe~do not exploit this
  %     representation in our own algorithm, but it helps to understand his final
  %     tree.}
  %~\end{figure}

%   The function is compiled to this decision tree: 

%   \begin{figure}[H]
%       \includegraphics[scale=0.7]{../images/dtree.png}
%       \Description{The final decision tree for merge} 
%       \caption{The final compiled decision tree for~\tt{merge}, right-to-left}
%   \end{figure}

%  The decision tree for~\tt{merge}, like the original function, tests values and
%  makes decisions. When presented with the values~\tt{xs} and~\tt{ys}, the tree
%  first tests ~\tt{xs} against its two known possible forms: the nullary list
%  constructor ~\tt{[]}, and an application of the~\it{cons} constructor~\tt{::}.
%  If ~\tt{xs} is equal to~\tt{[]}, the tree immediately returns~\tt{ys}. If
%  ~\tt{xs} is an application of~\tt{::}, the tree then tests~\tt{ys}
%  against~\tt{[]} and~\tt{::}, and it returns a value according to the result of
%  the match. Each time the tree goes down a~\tt{::} branch, it extracts the
%  arguments of the~\tt{::} for later use: these are~\tt{x}, ~\tt{y},~\tt{xr},
%  and~\tt{yr}, which are used in the~\tt{...} branch. This process of extracting
%  arguments generalizes to all value constructors with one or more arguments. 

  \begin{figure}
    \begin{center}
    \dcsyntax
    \end{center}
    \Description{The concrete syntax of D.}
    \caption{\D: Concrete syntax}
    \label{fig:dsyntax}
    \end{figure}

  % Here is the decision tree for~\tt{merge} in~\D. 

  %~\begin{figure}[H]
  %   \centering
  %     \it{Like the semantics, this figure is in progress.}
  %     \Description{The decision tree for merge, now in D} 
  %     \caption{The tree for~\tt{merge} in D looks similar to Maranget's.}
  %~\end{figure}

  In~\D, as in Maranget's trees, the~\it{test} node extracts all names from a
  value constructor at once for use in subtrees. The compiler is responsible for
  introducing the fresh names used in~\it{test}. The compiler alpha-renames all
  necessary terms before it translates an~\iffibf to a decision tree to ensure
  all names are unique. 

  In~\D, like in~\VMinus, expressions can~\it{fail}, meaning some of~\D's
  syntactic forms like~\it{try-let} and~\it{cmp} have an extra branch which is
  executed if the examined expression fails. 

    % Includes subsection
\rab{d semantics- show or omit?}
    % \label{dsemantics}
    % \dsemantics


    \subsection{The~\DTran\ algorithm:~\VMinus $\rightarrow$~\D}

    % To demonstrate that~\VMinus has a similarly-desirable cost model to pattern
    % matching, we~present an algorithm for compiling~\VMinus to a decision treeWe
    % choose the decision tree as a target for compilation for the simple reason
    % of its appealing cost model. A decision tree can be exponential in size but
    % never examines any word of the~\it{scrutinee}---the value being tested---%
    % more than once. This property is established by the compilation from~\VMinus
    % to \D, which ensuring that no~\it{test} node~\it{T} has any proper ancestor
    % \it{T'} such that~\it{T} and~\it{T'} both test the same location in memory.   

    % There are a few minor differences in the algorithm we~use and Maranget's: his
    % compilation algorithm is more complex than the one in this paper, and
    % involves an intermediate representation of occurrence vectors and clause
    % matrices which the algorithm we~present does not use. Maranget uses vectors
    % and matrices to express multiple simultaneous matches of values to patterns
    % as a single match of a vector with a matrix row. This allows him to run a
    % \it{specialization} pass that reduces the number of rows in the matrix,
    % ultimately leading to smaller trees. Because the trees produced by our
    % algorithm still have the linear-in-code-size property, we find them
    % acceptable for the current work. 

    % The algorithm runs during~\DTran, the transformation from~\VMinus to~\D. Its
    % domain, instead of a~\it{case} expression, is~\VMinus's~\iffibf.
    % \DTran~propagates a context~\ctx\ which maps each defined name to~\it{known} or
    % \it{unknown}. The context is used when determining the form of a guard.
    % in the compilation rules. 
       

    %~\algrenewcommand\algorithmicwhile{\bf{let }}
    %~\algrenewcommand\algorithmicdo{\bf{in }}
    %~\algrenewcommand\algorithmicend{\bf{end}}
    %~\newcommand\alet\algorithmicwhile
    %~\newcommand\ain\algorithmicdo
    %~\newcommand\aend\algorithmicend

    %~\newcommand\branches{\ensuremath{\mathit{branches}}}
    %~\newcommand\mg{\ensuremath{-}}


    When presented with an~\iffibf,~\DTran\~invokes~\Compile with context~\ctx.
    \Compile~first desugars the~\iffibf by expanding choice to multiple~\iffibf
    branches with a desugaring function~\ITran: 
    
    \itran{if\;~\dots~\dbar\; gs_{1};\;~\choiceg{gs_{2}}{gs_{3}};\; gs_{4}~\rightarrow e\;~\dbar~\dots~\;fi}
    
    ==
    
    ${if\;~\dots~\dbar\; gs_{1};\; gs_{2}~\rightarrow e~\;\dbar\; gs_{3};\; gs_{4}~\rightarrow e~\;\dbar~\dots~\;fi}$

    With the desugared~\iffibf,~\Compile\ then repeatedly chooses a guarded
    expression $G$ and applies one of the compilation rules below. The rules are
    applied in a nondeterministic order. 

    The algorithm terminates when it inserts a final~\it{match} node (with rule
    \textsc{Match}).  A~\it{match} node is inserted for a right-hand side expression~\expr when the list of
    guards preceding~\expr is empty or a list of assignments from names to
    unbound names. Termination of~\DTran\~is guaranteed because each recursive
    call passes a list of guarded expressions in which the number of guards is
    strictly smaller, so eventually the algorithm reaches a state in which the
    first unmatched branch is all trivially-satisfied guards.


    If~\Compile\ cannot choose a $g$ of
    one of the valid forms, it halts with an error. This
    can happen when no $g$ is currently solvable in the context, as determined
    by the same algorithm that~\VMinus uses to pick a guard to solve, or when
    the program would be forced to unify incompatible values, such as any value
    with a closure. 

    \rab{I have these rules with english descriptions, though that increases length}.

    
\showvjudgement{Compile}{\compilebig}

      \subsection{Reduction Strategies}

      In our implementation, we apply the~\textsc{Test} rule before other rules,
      having found in our experiments that this heuristic leads to the smallest trees. A~%
      risk of inserting a \it{test} node before a \it{let-unless} node is that
      if there are many shared names across branches, a \it{test} node
      will introduce those names in \it{let-unless} nodes in each
      subtree.
      It is possible to insert \it{let-unless} nodes before \it{test}
      nodes, in which case those names will be bound
      only once.  The risk of the \it{let}-first strategy
      is that if there are many names used in only one subtree of a \it{test},
      these names will be introduced to \it{all} branches.  The extra
      bindings may harm
      performance by bloating an environment in an interpreter or thrashing the
      icache if the tree further is compiled to machine code. 

      In the implementation,~\DTran\~first introduces all the names under the
      all existential $\exists$'s to a context which determines if a name is
      \it{known} or~\it{unknown}, applying the~\textsc{Exists} rule is applied
      all at the start. At the start of compilation, each name introduced by
      $\exists$ is~\it{unknown} in a context~\ctx. Since all names in the
      program are unique at this stage, there are no clashes. 

      \subsection{Full Big-step rules for \Compile, with no descriptions}

        The judgement form for compilation is \textsc{Compile}: 

        \showvjudgement{Compile}{\compilebig}

The rules are nondeterministic: the structure of the final result~\expr\ depends
on the order in which rules are applied by the compiler.  

\rawcompiler


    \section{Implementations}

    We~have placed implementations of~\VMinus and~\D at
    \url{https://github.com/rogerburtonpatel/vml}. The implementations are
    complete, from parsers to evaluation to unparsers. In the same repository
    lives the~\tt{dtran} program, which translates from~\VMinus to~\D. With the
    implementations, we~include test cases for
    evaluation of~\VMinus, evaluation
    of~\D, and the translation between the two. 
    In~each of these test cases, the translation preserves semantics. 


    \section{Related and Future Work}

    This paper builds on Augustsson et al.'s Verse Calculus~\citep{verse} and
    decision trees~\citep{maranget}. Augustsson et al. give the formal rewrite
    semantics for the Verse Calculus; Maranget gives an elegant formalism of
    decision trees and a translation algorithm from patterns to decision trees.
    We~attempted to imitate the behavior of the rewrite semantics of~\VC in the
    big-step semantics of~\VMinus by manually rewriting terms in~\VC and
    creating rules that would imitate the ultimate result of term-rewriting.
    Proving the equivalence between the two semantics is the subject of future
    workWe~chose a big-step semantics because it is the style of semantics we~are
    most comfortable with; writing the formalisms this way helped me write the
    code. Using a rewrite semantics instead would more closely relate~\VMinus
    and~\VC, and is a likely future project. Maranget's formalism was the
    foundation off of which we~built~\D. 
    
    Extensions to pattern matching, and how they appeal to language designers,
    find an excellent example in~\citet{guardproposal}. The authors describe
    pattern guards and transformational patterns (another extension to pattern
    matching), both of which allow a Haskell programmer to write more concise
    code using pattern matching. Or-patterns are documented in the OCaml
    Language Reference Manual~\citep{ocaml}.
    
    \citet{augustsson1985compiling}~ gives a foundation in compiling pattern
    matching. \citet{scottramsey} have a crisp example of a match-compilation
    algorithm (pattern matching to decision trees). Scott and Ramsey's algorithm
    structurally inspired mine, and studying the source code from the paper
    aided our implementation. 

    %\subsubsection{Desirable theorems about translations}
    For future work, our top priority is to prove that \DTran\ preserves
    semantics.
    %Solid proofs~will strengthen our
    %argument that~\VMinus is not only a viable alternative to pattern
    %matching syntactically; it is also formally equivalent. 
    Next, we~plan to prove that
    %Two theorems would strengthen the viability of~\VMinus, and are
    %essential targets for future work:~\bf{Proving~
    \VMinus~is deterministic, and that \bf{the big-step semantics of~\VMinus~is
    consistent with the published semantics of~\VC.} As the authors of the Verse
    paper proved that the rewrite semantics of Verse is skew-confluent, we~plan
    to prove that big-step semantics of~\VMinus is deterministic, despite the
    nondeterminism of choosing a guard. Second, \VMinus~is designed to be
    Verse-like, and formalizing the relationship between the two would
    strengthen \VMinus's~viability as a language that compromises between
    pattern matching and equations. 

    Finally, we~would like to explore
    exhaustiveness analysis of~\VMinus.
    Exhaustiveness analysis
    can warn programmers of a missing or extraneous alternative in a~\it{case} expression.
    Owing to its significantly more flexible structure,
    however,~\iffibf may prove trickier to analyze.
    
    
    % \section{Discussion: The design of \VMinus}        
    % \label{futurework}
    
    %~\subsubsection{\VMinus and~\PPlus, side by side}
    % We~compare~\VMinus with~\PPlus as an exercise in comparing equations with
    % pattern matching. They certainly look similar, which hints that~\VMinus
    % might be as expressive as pattern matching with the three extensions. 
    % Proving this claim is the topic of Section~\ref{pplustovminus}. 
    
    % When might a programmer prefer~\VMinus over~\PPlus, or vice versa? After
    % programming in both, we~have come up with three empirical observations: 
    
    %~\begin{enumerate}
    %   \item~\it{The scrutinee:} When there is no obvious single scrutinee,~\VMinus is more succinct.
    %   When there is a scrutinee,~\PPlus is more succinct. 
    %   \item~\it{Binding and decision-making:} Binding and decision-making are
    %   joined in a single construct in~\VMinus:~\tt{=}.~\PPlus needs different
    %   kinds of syntax, like~\tt{<-}, to express different kinds of binding. And
    %   a programmer wanting for~\tt{let} in~\PPlus will never feel this need in
    %   \VMinus:~\tt{=} subsumes that, too. 
    %   \item~\it{Names:} In~\VMinus, names are explicitly introduced. In
    %   practice, this helps prevent a common mistake in pattern matching in which
    %   a programmer attempts to match a value $v$ on an in-scope name $x$,
    %   expecting the match to succeed iff $x$ evaluates to $v$ at runtime, only
    %   to see the match always succeed. 
    %~\end{enumerate}     
    
    %~\rab{How do you feel about examples for each of these points here? This
    % section is very long, but they help drive these points home.}
    
    
    % An example in~\VC illustrates this precise
    % scenario: 
    
    % FIRST PAPER EXAMPLE 
    
    % IMPORTANTLY,~\VMinus REALLY ISN'T A FUNCTIONAL LOGIC PROGRAMMING LANGUAGE ANYMORE. 
    
    %  Doesn't backtrack
    
    %  No multiple results 
    
    %  Doesn't evaluate functions backwards, have top-level patterns like verse, the list goes on 
    
    % If you're well-versed in functional logic programming, may perceive that
    % imposing such restrictions on choice and names effectively strips away
    % much of Verse's essence as a functional logic programming language. With
    % these constraints enforced, there can be no backtracking, multiple
    % results, backward function evaluation, or top-level patterns, among
    % other classic functional logic programming features. But do not fear:
    % our intent is not to recklessly strip Verse of its meticulously crafted
    % core tenets. Instead, our aim is to extract a select few—namely, its
    % equations, existentials, and nondeterministic evaluation order—and
    % juxtapose them with pattern matching.    
    
    
    %~\section{Addressing how~\PPlus handles unusual pattern combinations}
    %~\label{ppweird}
    %     \PPlus admits of strange-looking patterns: consider~\tt{Cons (when true)
    %     zs}. But these should not be alarming, because such syntactic forms reduce
    %     to normal forms by (direct) application of algebraic laws: 
    
    %     \begin{align}
    %       K (\whenexpr)~\;p2\;~\dots &=== K~\;\_~\;p2\;~\dots,~\whenexpr~\\
    %       K (\whenexpr,~\;p2\;)~\;p3\;~\dots  &=== K~\;p2\;~\;p3\;~\dots,~\whenexpr~\\
    %       K (p1\;,~\whenexpr)~\;p3\;~\dots  &=== K~\;p1\;~\;p3\;~\dots,~\whenexpr~\\
    %       K (\whenexpr~\pbar p2\;)~\;p3\;~\dots &=== (K~\;\_~\;p3\;~\dots,~\whenexpr)~\pbar (K~\;p2\;~\;p3\;~\dots)~\\
    %       K (p1~\pbar~\whenexpr)~\;p3\;~\dots &=== (K~\;p2\;~\;p3\;~\dots)~\pbar (K~\;\_~\;p3\;~\dots,~\whenexpr)  \\
    %       \whenexpr~\leftarrow e &===~\;\_ <- e,~\whenexpr
    %     \end{align}   
        
    %     Repeatedly applying these laws until the program reaches a fixed point
    %     normalizes placements of~\it{when}. Laws (2) and (3) work because~\PPlus
    %     has no side effects and the laws assume all names are unique (the compiler
    %     takes care of this), so changing the order in which patterns match has no
    %     effect on semantics.         
    

%    \subsection{The design of~\VMinus}

    % This section is aimed at audience familiar with \VC; in particular, with the
    % concept of \it{names as values}.
    
    % \subsection{Forms in \VC and \VMinus}
    
    % When designing~\VMinus, we~wanted the language to capture the expressiveness
    % of~\VC's equations while retaining a similar decision-making construct to
    % pattern matching. In~pattern matching, the only decision-making
    % construct is~\it{case}; other forms of decision-making like~\it{if}
    % desugar to it. In~\VMinus, the only decision-making construct is~\iffibf.
    % This design differs from~\VC, which features numerous ways to make decisions by
    % combining~\bf{one}, \bf{all}, equations, intermediate expressions,
    % and~choiceWe~did not want multiple results in~\VMinus, so we~eliminated
    % anything that looked like~\bf{all}, and we~combined all of the above
    % constructs into the singular form~\iffibf. Like~\it{case}, there is only one
    % way to use~\iffibf- unlike \bf{one}, equations, intermediate expressions,
    % and~choice in~\VC. By restricting all decision-making to \iffibf,
    % programmers cannot ``misuse'' any of the forms in ways that might lead to
    % problematic computations, such as multiple results by returning choice as an
    % expression. Furthermore, as we've shown in the examples, the way in which programmers use~\iffibf mirrors
    % the way they use~\it{case}. 
    
    % \subsection{Choice in \VMinus vs. \VC}

    % \VC's choice operator is often the culprit behind both
    % backtracking and multiple results, which tempted me to remove choice
    % from~\VMinus altogether. However, we~want to harness the expressive potential
    % of choice, particularly when paired with~\VC's~\bf{one} keyword. When
    % combined with choice, \bf{one} elegantly signifies ``proceed if any branch of
    % the choice succeeds.'' 
    
    % To this end, in~\VMinus, choice is permitted with several modifications:
    
    %     \begin{enumerate}
    %     \item Choice may only appear as a condition or 'guard', not as a result
    %     or the right-hand side of a binding.
    %     \item If any branch of the choice succeeds, the choice succeeds,
    %     producing any bindings found in that branch. The program examines the
    %     branches in a left-to-right order.
    %     \item The existential $\exists$ may not appear under choice.
    %     \end{enumerate}
        
    %     We~introduce one more crucial modification to~\VC: a name in
    %     \VMinus is an~\it{expression}, not a~\it{value}. This modification,
    %     coupled with our adjustments to choice, eradicates backtracking. Our~rationale is straightforward: if~an expression returns a
    %     name, and if the program later imposes a new constraint on that
    %     name, it may necessitate the reevaluation of the earlier expression---a
    %     scenario we~insist on avoiding. 



        %~\subsubsection{Type-agnostic decision-making: the $\alpha$}
        %~\label{alphas}

        % The three languages look similar: they each have value constructors and
        % a 'decision-making construct' to deal with constructed data. In~\PPlus, the
        % construct is pattern-matching; in~\VMinus, it is the guarded expression; in~\D,
        % it is the decision tree. 

        % Because of this, it might be possible to make all three languages
        %~\it{higher order} in right-hand sides; that is, to parameterize the
        % expressions that occur~\it{after} all decision-making. Imagine an
        % abstract expression $\alpha$ that appears on the right-hand side of a
        %~\it{case} branch, the right-hand side of a guarded expression, or in a
        %~\it{match} node. The abstract syntax of the new~\it{case},~\iffibf,
        % and decision tree might look like this: 
        %~\begin{center}
        %     \begin{bnf}
        %         $\VMinus_{\alpha}$ :~\VMinus with $\alpha$ ::=
        %         $\mathit{if}\;~\mathit{[}\; g\;~\bracketed{[] g}\;~\mathit{]}\;~\mathit{fi}$ : if-fI~with $\alpha$
        %         ;;
        %         $G_{\alpha}$ : Guarded Expressions with $\alpha$ ::=
        %         $[\vexists{\bracketed{x}}]~\bracketed{g}~\boldsymbol{\rightarrow}\alpha$ : 
        %         ;;
        %         $\PPlus_{\alpha}$ :~\PPlus with $\alpha$::=
        %         $\tt{case}\;~\expr\;~\bracketed{p\;~\ttrightarrow\;~\alpha}$ :~\it{case} expression with $\alpha$ 
        %         ;;
        %         $t_{\alpha}$ : Decision Trees with $\alpha$ ::= 
        %         |~\dots : other forms of tree 
        %         | $\alpha$ : match node 
        %     \end{bnf}
        %~\end{center}

        % Why would one want to do this? Well, recall that~\VMinus had to be
        % stripped of multiple results and other~\VC-like constructs in order to
        % retain its desirable efficiency properties. $\alpha$ lets~\VMinus and
        % the other languages do efficient decision-making without worrying about
        % the form of the result. If the result is a complex multi-value or a
        % computation that involves backtracking, it will be agnostic of the
        % decision-making. $\alpha$ makes right-hand sides polymorphic and
        % abstract, so a programmer could potentially insert expressions from~\VC
        % in their place and know that the decision-making before the $\alpha$
        % will still be efficient. This would allow for fuller interoperability
        % between~\VC and~\VMinus. Section~\ref{vminusandvc} further describes
        % why bridging the gap between the two languages might be a worthwhile
        % exercise. 

        % However, the language designer must take special care to ensure no
        % $\alpha$ finds its way into the decision-making itself, or the whole
        % idea falls apart. We~are developing an implementation that enforces this
        % invariant, and may include it in a future publication. 
        
        % Other alpha material: 
% The decision-making construct that gets us there. Whether
% it's a single value (ML-style) a sequence of values (Verse-style), or
% even something else, the $\alpha$ represents~\it{any} ultimate result of
% "making a decision," and it's the ways in which we make decisions that
% we truly care about examining. By making the return result both
% polymorphic and abstract, we eschew the need to worry about its type and
% compatibility with other results of otherwise-equivalent trees. 

% An expression in core Verse evaluates to produce possibly-empty sequence of
% values. In~\VMinus, values depend on the form of abstract expression $\alpha.$
% If $\alpha$ is a Verse-like expression,~\valpha\ will be a value sequence. If it
% is an ML-like expression, it will be a single value. 
            
%         A guarded expression evaluates to produce a~\bf{result}. A result is either
%         a metavalue~\valpha\ or reject. 
            
%         \[\it{r}\;~\rm{::=}\;~\vartheta\;~\vbar~\;~\reject~\]
            
%         \showvjudgement{Eval-Alpha}{\veval{\alpha}{\valpha}}

% Of note in both~\VMinus and~\D is that the 'decision-making construct'
% is annotated with an $\alpha$. This annotation gives us type flexibility on the
% right-hand side of the~\it{terminating} case for each construct
% (\tt{$\rightarrow$ exp} in~\VMinus and the match node in~\D.) 


        %~\subsubsection{Or-patterns and pattern guards}
        %~\label{pplusindependently}
        %~\PPlus has side conditions, guards, and or-patterns. No major
        % functional language has all three of these extensions. Back
        % Section~\ref{extensions}'s examples, we~had to switch from OCaml to
        % Haskell to use guards, and back to OCaml for or-patterns. The two
        % extensions are mutually exclusive in Haskell, OCaml, Scala,
        % Erlang/Elixir, Rust, F\#, Agda.~\citep{haskell, ocaml, scala, erlang,
        % elixir, rust, fsharp, agda}

        % Why is this the case? In our research, we~have yet to encounter a
        % substantial justification for thisWe~have several theories: one,
        % reengineering the Haskell parser to integrate or-patterns into the
        % language may be considered too great an effort; two, the lesser
        % popularity of functional programming in comparison to other paradigms
        % has meant there are not enough voices in any one language's community
        % claiming that theirs needs all known extensions to pattern matching;
        % three, Haskell, the only language with guards, is lazy, and lazy
        % semantics may not operate well with or-patterns. Future work may explore
        % if a lazy semantics of~\PPlus is possible; such a study may answer this
        % question. 

        % \subsubsection{Using~\VMinus to inform programming in Verse}
        % \label{vminusandvc}
        
        % At ICFP last year, Tim Sweeney said that he wanted Verse to be an
        % accessible programming language to write a scalable, collaborative
        % metaverse~\citep{timtalk}. Can~\VMinus be aid in this goal? we~can imagine
        % two ways in which it might:
        
        %~\begin{enumerate}
        %     \item~\VMinus could be  a tool to help ease programmers who are
        %     more familiar with pattern matching into the realm of functional
        %     logic programming with equations. 
        %     \item Programs written in Verse using ideas from~\VMinus might have
        %     a friendlier cost model (depending on the compiler)
        %~\end{enumerate}
        
        % To point~1,~\VMinus sits both syntactically and semantically in between
        %~\PPlus and~\VC, which might help a new programmer to Verse bridge the
        % conceptual gap between pattern matching and equations. Also,~\PtoVTran,
        % the $\PPlus~\rightarrow~\VMinus$ translation, could help a programmer 
        % who wishes to write code using pattern matching see how their ideas 
        % can be expressed in Verse. 

        % To point~2, 
        
        % \DTran\ and the proof that~\DTran\ preserves semantics help show that
        % certain computations that use equations for decision-making can be
        % compiled to efficient code. A future project could be to extend~\VMinus
        % to include~\it{all} of~\VC, and use~\DTran to eliminates as much
        % backtracking at runtime as possible, falling back to the VC's fully
        % general evaluation mechanism only when necessary. Our hope is that, using
        % these ideas, both the Verse programmer and language designer might make
        % any discovery that allows them to increase the efficiency of full-Verse
        % programs. 

    \section{Conclusion}

    We~have introduced \VMinus, a~language that makes decisions using equations.
    By~example, we have shown that programs written in~\VMinus can have the same
    desirable properties as equivalent programs written using pattern matching.
    And~to~show that equations can be compiled to efficient decision trees, we
    have introduced~\D and~\DTran. In doing so, we~have demonstrated that
    programming with equations is a promising alternative to pattern matching.
    %%  \D~and~\DTran~serve to demonstrate that~\VMinus is a viable alternative to
    %%  pattern matching on the grounds that they have equivalent cost models:
    %%  pattern matching can be compiled to a decision tree,
    %%  which~\citet{macqueen1985tree} built the foundation for and~\citet{maranget}
    %%  expanded on.

    We~have also fully implemented the languages. They exist for use and
    experimentation: they are syntactically simple and have conceptually
    accessible operational semanticsWe~hope that programmers will explore and
    develop their own opinions of these languages, which are publically
    available at~\url{https://github.com/rogerburtonpatel/vml}. 

    % Finally, and in particular with~\VMinus, we~hope to provide a stepping
    % stone between pattern matching and equations that a new programmer to
    % Verse will find illuminating. 
    
    % hope to have paved a small
    % segment of the path that the curious programmer or language enthusiast who
    % wishes to better understand Verse will take. Be they transitioning to
    % equations from pattern matching to equations or curious about how those
    % equations might be compilable to decision trees, we~hope they find the
    % languages, and this document, illuminating. 
    



\end{document}
\endinput
%%
